\section{“Teach Them What They Should Know for Salvation”}

The most probative interpretation of Adele's reported commission is her response. By 1871, before the fire tradition had become the dominant public image, Starr encountered a visible catechetical project. Pernin in 1874 likewise describes land, buildings, instruction, collaborators, and pilgrimage. The mission is therefore not a moral appended by a later shrine to an otherwise empty wonder; it appears in the earliest located printed evidence.

\subsection{From the road to the household}

Later accounts describe Adele traveling from home to home, sometimes over considerable distances, to teach children and persuade families to support religious instruction. The work required access, trust, food, lodging, books, space, and cooperation. “Gather” thus took practical forms before it took architectural form.

Household visitation resists two reductions of catechesis. First, children are not isolated minds into which facts are deposited; their religious life is embedded in family and community. Second, faith cannot be reduced to domestic inheritance without ecclesial teaching and sacramental incorporation. Adele's apostolate moves between home and Church.

\subsection{The early institutional fruit}

Starr's 1871 narrative records a free school and a fundraising journey for a school, chapel, and home. Pernin's 1874 appendix describes the site and the women who assisted the work. These institutions matter historically because they existed close enough to the reported event to establish sustained enactment. They matter theologically because charisms are tested in service and endurance rather than intensity of recollection.

Fruit is evidence, not a mechanical proof. A useful institution can grow around a mistaken origin story, and genuine grace can pass through fragile institutions. The 2010 discernment therefore considers fruit together with doctrine, character, testimony, and the absence of evident manipulation. The school cannot prove the apparition alone; it can show what Adele understood the apparition to require.

\subsection{The fourfold pedagogy}

The received command contains a compact program of Christian initiation:

\begin{fourstable}{Element}{What is taught}{What it resists}{Mature form}
What is needed for salvation & The saving work of the Triune God, faith, conversion, commandments, grace, prayer, and Christian life & Vague uplift with no account of Christ or response & A coherent proclamation ordered to communion with God, proportioned to the learner\\
The catechism & An ecclesially received synthesis of doctrine & Private invention and disconnected religious impressions & Teaching that uses the Church's current catechetical norms and sound texts\\
The Sign of the Cross & A bodily confession of Trinity and Paschal mystery & A split between ideas and prayer & Reverent practice whose words are understood and lived\\
Approach to the sacraments & Preparation for encounter with Christ in the Church, especially Penance and Eucharist in this story & Treating sacraments as rewards, magic, or social milestones & Faith, repentance, proper disposition, mystagogy, and continuing participation\\
\end{fourstable}

The list is not exhaustive, but its integration is notable. Doctrine without prayer becomes abstraction. Gesture without doctrine becomes habit emptied of reference. Sacramental access without conversion becomes formalism. Zeal without ecclesial content makes the catechist the measure of truth. Champion holds the elements together.

\subsection{Salvation is Christ, not an examination syllabus}

The phrase “what they should know for salvation” could be misheard as though heaven supplied a minimum-answer test. Catholic catechesis is intellectual because truth matters, but salvation is not information alone. It is the gracious communion with God opened by Christ, received in faith, expressed in love, and nourished sacramentally. Necessary knowledge is relational and practical: Who is God? What has Christ done? What does conversion require? How does one pray and live as a member of his Body?

The Church's catechetical directories and tradition describe catechesis as more than classroom instruction. It initiates into knowledge of faith, liturgical life, moral formation, prayer, community, and mission. Champion's early vocabulary is simpler, but the trajectory is compatible: children are led from elementary gesture and teaching toward sacramental belonging and a life shaped by the Gospel.

\subsection{Parents, pastors, and the catechist}

Parents possess a primary responsibility for their children's Christian education. Pastors have responsibilities arising from ecclesial office. Lay catechists participate by baptism, competence, commissioning, and service. Adele's story should not be used to pit these vocations against one another. Her work responds to a gap by engaging families and orienting children toward the sacraments administered in the Church.

Nor does an alleged private command dispense with formation. Contemporary catechists should know Scripture and doctrine, understand human development, learn sound pedagogy, honor disability and cultural difference, and follow safeguarding norms. Zeal can begin a journey; it cannot substitute for competence where children's welfare is entrusted to someone.

\subsection{The Sign of the Cross as miniature catechesis}

The Sign of the Cross condenses Champion's theological center. Its words confess Father, Son, and Holy Spirit. Its gesture traces the instrument of Christ's self-giving. It is received in baptismal culture and repeated at prayer and liturgy. Teaching a child to make it well can therefore unite memory, body, doctrine, and belonging.

The sign is not a protective charm. Christian use does not promise immunity from fire, illness, or grief. It marks the person with the Paschal mystery: the victory of love is revealed through, not around, the Cross.

\subsection{Adele's poverty and the economy of collaboration}

Fundraising appears in the earliest printed witness. This prevents an over-spiritualized portrait. A free school still needs land, fuel, food, maintenance, teachers, and shelter. Dependence on alms can manifest trust, but it can also expose a work to instability. Adele's mission becomes sustainable through benefactors, companions, clergy, families, and eventually institutional stewardship.

The assurance of help is thus not narrated as freedom from means. Providence arrives through human generosity and disciplined labor. That pattern supplies a criterion for authentic reception today: prayer for help should increase responsibility, transparent stewardship, and collaboration, not evade them.

\subsection{Women as agents of transmission}

The narrative's central human actor is a lay immigrant woman; its earliest public work is shared with other women; many of its first learners were children beyond the easy reach of institutions. This does not create a theory of ecclesial office from an apparition. It does reveal something already grounded in baptism: women are indispensable bearers, teachers, organizers, and witnesses of the faith.

Adele's authority grows by service rather than self-advertisement. Her model is neither silent passivity nor domination clothed as charisma. It is durable, relational agency ordered to another's good.

\interpretivekey{The apparition interpreted by its fruit}{Champion's catechetical fruit does not relieve the historian of source criticism. It tells the theologian where to look for the message's meaning. The extraordinary claim is ordered toward the ordinary means by which Christ forms a people: proclamation, patient teaching, repentance, sacrament, prayer, family, and works of mercy.}
